% assets/w07-atan2-alpha.svg — atan2 가 왜 그 각인가
%
%   bash _tools/tikz2svg.sh assets/src/w07-atan2-alpha.tex assets/w07-atan2-alpha.svg
%
% 요점 둘
%   (1) 배 자리의 두 성분은 서로 직각이다. 그래서 둘은 직각삼각형의 밑변과
%       높이이고, alpha 는 밑변에서 빗변까지의 각이다. tan(alpha) = 높이/밑변.
%   (2) 비 하나만으로는 각이 정해지지 않는다. 같은 비를 주는 점이 반대쪽에
%       하나 더 있고, 그 둘의 침로는 180도 어긋난다. 그래서 atan 이 아니라 atan2.
%
% 스케일: 1 m/s = 18 mm  (1-5절 기하 그림의 속도 삼각형과 같다)
%   rdot = -0.859 -> -15.5 mm,  r*thetadot = 1.229 -> 22.1 mm,  v_d = 1.5 -> 27 mm
%   alpha = atan2(1.229, -0.859) = 124.9 deg
%   반대쪽 점 (0.859, -1.229) 의 각 = -55.1 deg,  두 점의 비는 똑같이 -1.43

\documentclass[border=6pt]{standalone}
% capstone-style.tex — 이 볼트의 TikZ 그림이 공유하는 스타일
%
% 쓰는 법:  % capstone-style.tex — 이 볼트의 TikZ 그림이 공유하는 스타일
%
% 쓰는 법:  % capstone-style.tex — 이 볼트의 TikZ 그림이 공유하는 스타일
%
% 쓰는 법:  \input{capstone-style}   (assets/src/ 안에서)
% 빌드:     bash _tools/tikz2svg.sh assets/src/w02-node-topic.tex assets/w02-node-topic.svg
%
% 왜 TikZ 인가
%   손으로 SVG 좌표를 쓰면 화살촉·선 굵기·글자 크기가 그림마다 미묘하게 달라진다.
%   그것이 "자료가 어설퍼 보인다"의 정체다. 스타일을 한 번만 정하고 상대좌표로
%   배치하면 좌표를 고쳐도 그림이 무너지지 않는다.
%   수식은 본문과 같은 TeX 글꼴로 나온다.
%
% 한글
%   ko.TeX(DVI 경로)를 쓴다. XeLaTeX 은 TikZ 그래픽을 PDF special 로 내보내는데
%   dvisvgm 이 그것을 받으려면 Ghostscript < 10.01 이 필요하다. 이 PC 에는
%   10.07 이 깔려 있어 그 경로가 막힌다. latex -> DVI -> dvisvgm 이 유일하게
%   동작하는 조합이다.

\usepackage[hangul]{kotex}
\usepackage{tikz}
\usepackage{amsmath}
\usepackage{xcolor}
\usetikzlibrary{arrows.meta, positioning, calc, fit, backgrounds, decorations.pathmorphing}

% ---- 색 --------------------------------------------------------------------
% 강조는 하나만. 나머지는 검정. 흑백 인쇄에서도 읽혀야 한다.
\definecolor{ink}{HTML}{1A1A1A}
\definecolor{soft}{HTML}{666666}
\definecolor{accent}{HTML}{7C3AED}
\definecolor{warn}{HTML}{B91C1C}

% 계층 색 — Simulink 모델의 블록 색과 같은 규칙을 쓴다
\definecolor{guid}{HTML}{1D4ED8}    % 유도
\definecolor{ctrl}{HTML}{EA580C}    % 제어
\definecolor{thr}{HTML}{D97706}     % 추진기
\definecolor{plant}{HTML}{16A34A}   % 운동모델
\definecolor{nav}{HTML}{7C3AED}     % 항법 · 신호처리
\definecolor{pathc}{HTML}{0D9488}   % 계획 경로 · 항적
\definecolor{errc}{HTML}{C2410C}    % 오차
\definecolor{flow}{HTML}{0369A1}    % 조류 · 바람

% ---- 선과 화살촉 -----------------------------------------------------------
\tikzset{
  sig/.style   = {ink, line width=0.5pt, -{Stealth[length=4pt, width=3pt]}},
  wire/.style  = {ink, line width=0.5pt},
  fb/.style    = {ink, line width=0.5pt, densely dashed,
                  -{Stealth[length=4pt, width=3pt]}},
  asig/.style  = {accent, line width=0.5pt, densely dashed,
                  -{Stealth[length=4pt, width=3pt]}},
  % 블록 — 흰 바탕, 직각, 검은 테두리. 색은 테두리에만 준다
  blk/.style   = {draw=ink, line width=0.55pt, fill=white,
                  minimum width=18mm, minimum height=8.5mm, inner sep=2pt, font=\small},
  gblk/.style  = {blk, draw=guid},
  cblk/.style  = {blk, draw=ctrl},
  tblk/.style  = {blk, draw=thr},
  pblk/.style  = {blk, draw=plant},
  nblk/.style  = {blk, draw=nav},
  % 합산점 · 분기점
  sum/.style   = {draw=ink, line width=0.55pt, fill=white, circle,
                  inner sep=0pt, minimum size=4.6mm, font=\scriptsize},
  dot/.style   = {ink, fill=ink, circle, inner sep=0pt, minimum size=1.5mm},
  % 글자
  sname/.style = {font=\small, inner sep=1.5pt},
  note/.style  = {font=\scriptsize, soft, inner sep=1.5pt},
  wnote/.style = {font=\scriptsize, warn, inner sep=1.5pt},
  anote/.style = {font=\scriptsize, accent, inner sep=1.5pt},
  axis/.style  = {ink, line width=0.5pt},
  tick/.style  = {font=\scriptsize, soft},
  % 묶음 상자 — 서브시스템 표시
  grp/.style   = {draw=soft, line width=0.4pt, densely dotted, rounded corners=1mm,
                  inner sep=3mm},
}
   (assets/src/ 안에서)
% 빌드:     bash _tools/tikz2svg.sh assets/src/w02-node-topic.tex assets/w02-node-topic.svg
%
% 왜 TikZ 인가
%   손으로 SVG 좌표를 쓰면 화살촉·선 굵기·글자 크기가 그림마다 미묘하게 달라진다.
%   그것이 "자료가 어설퍼 보인다"의 정체다. 스타일을 한 번만 정하고 상대좌표로
%   배치하면 좌표를 고쳐도 그림이 무너지지 않는다.
%   수식은 본문과 같은 TeX 글꼴로 나온다.
%
% 한글
%   ko.TeX(DVI 경로)를 쓴다. XeLaTeX 은 TikZ 그래픽을 PDF special 로 내보내는데
%   dvisvgm 이 그것을 받으려면 Ghostscript < 10.01 이 필요하다. 이 PC 에는
%   10.07 이 깔려 있어 그 경로가 막힌다. latex -> DVI -> dvisvgm 이 유일하게
%   동작하는 조합이다.

\usepackage[hangul]{kotex}
\usepackage{tikz}
\usepackage{amsmath}
\usepackage{xcolor}
\usetikzlibrary{arrows.meta, positioning, calc, fit, backgrounds, decorations.pathmorphing}

% ---- 색 --------------------------------------------------------------------
% 강조는 하나만. 나머지는 검정. 흑백 인쇄에서도 읽혀야 한다.
\definecolor{ink}{HTML}{1A1A1A}
\definecolor{soft}{HTML}{666666}
\definecolor{accent}{HTML}{7C3AED}
\definecolor{warn}{HTML}{B91C1C}

% 계층 색 — Simulink 모델의 블록 색과 같은 규칙을 쓴다
\definecolor{guid}{HTML}{1D4ED8}    % 유도
\definecolor{ctrl}{HTML}{EA580C}    % 제어
\definecolor{thr}{HTML}{D97706}     % 추진기
\definecolor{plant}{HTML}{16A34A}   % 운동모델
\definecolor{nav}{HTML}{7C3AED}     % 항법 · 신호처리
\definecolor{pathc}{HTML}{0D9488}   % 계획 경로 · 항적
\definecolor{errc}{HTML}{C2410C}    % 오차
\definecolor{flow}{HTML}{0369A1}    % 조류 · 바람

% ---- 선과 화살촉 -----------------------------------------------------------
\tikzset{
  sig/.style   = {ink, line width=0.5pt, -{Stealth[length=4pt, width=3pt]}},
  wire/.style  = {ink, line width=0.5pt},
  fb/.style    = {ink, line width=0.5pt, densely dashed,
                  -{Stealth[length=4pt, width=3pt]}},
  asig/.style  = {accent, line width=0.5pt, densely dashed,
                  -{Stealth[length=4pt, width=3pt]}},
  % 블록 — 흰 바탕, 직각, 검은 테두리. 색은 테두리에만 준다
  blk/.style   = {draw=ink, line width=0.55pt, fill=white,
                  minimum width=18mm, minimum height=8.5mm, inner sep=2pt, font=\small},
  gblk/.style  = {blk, draw=guid},
  cblk/.style  = {blk, draw=ctrl},
  tblk/.style  = {blk, draw=thr},
  pblk/.style  = {blk, draw=plant},
  nblk/.style  = {blk, draw=nav},
  % 합산점 · 분기점
  sum/.style   = {draw=ink, line width=0.55pt, fill=white, circle,
                  inner sep=0pt, minimum size=4.6mm, font=\scriptsize},
  dot/.style   = {ink, fill=ink, circle, inner sep=0pt, minimum size=1.5mm},
  % 글자
  sname/.style = {font=\small, inner sep=1.5pt},
  note/.style  = {font=\scriptsize, soft, inner sep=1.5pt},
  wnote/.style = {font=\scriptsize, warn, inner sep=1.5pt},
  anote/.style = {font=\scriptsize, accent, inner sep=1.5pt},
  axis/.style  = {ink, line width=0.5pt},
  tick/.style  = {font=\scriptsize, soft},
  % 묶음 상자 — 서브시스템 표시
  grp/.style   = {draw=soft, line width=0.4pt, densely dotted, rounded corners=1mm,
                  inner sep=3mm},
}
   (assets/src/ 안에서)
% 빌드:     bash _tools/tikz2svg.sh assets/src/w02-node-topic.tex assets/w02-node-topic.svg
%
% 왜 TikZ 인가
%   손으로 SVG 좌표를 쓰면 화살촉·선 굵기·글자 크기가 그림마다 미묘하게 달라진다.
%   그것이 "자료가 어설퍼 보인다"의 정체다. 스타일을 한 번만 정하고 상대좌표로
%   배치하면 좌표를 고쳐도 그림이 무너지지 않는다.
%   수식은 본문과 같은 TeX 글꼴로 나온다.
%
% 한글
%   ko.TeX(DVI 경로)를 쓴다. XeLaTeX 은 TikZ 그래픽을 PDF special 로 내보내는데
%   dvisvgm 이 그것을 받으려면 Ghostscript < 10.01 이 필요하다. 이 PC 에는
%   10.07 이 깔려 있어 그 경로가 막힌다. latex -> DVI -> dvisvgm 이 유일하게
%   동작하는 조합이다.

\usepackage[hangul]{kotex}
\usepackage{tikz}
\usepackage{amsmath}
\usepackage{xcolor}
\usetikzlibrary{arrows.meta, positioning, calc, fit, backgrounds, decorations.pathmorphing}

% ---- 색 --------------------------------------------------------------------
% 강조는 하나만. 나머지는 검정. 흑백 인쇄에서도 읽혀야 한다.
\definecolor{ink}{HTML}{1A1A1A}
\definecolor{soft}{HTML}{666666}
\definecolor{accent}{HTML}{7C3AED}
\definecolor{warn}{HTML}{B91C1C}

% 계층 색 — Simulink 모델의 블록 색과 같은 규칙을 쓴다
\definecolor{guid}{HTML}{1D4ED8}    % 유도
\definecolor{ctrl}{HTML}{EA580C}    % 제어
\definecolor{thr}{HTML}{D97706}     % 추진기
\definecolor{plant}{HTML}{16A34A}   % 운동모델
\definecolor{nav}{HTML}{7C3AED}     % 항법 · 신호처리
\definecolor{pathc}{HTML}{0D9488}   % 계획 경로 · 항적
\definecolor{errc}{HTML}{C2410C}    % 오차
\definecolor{flow}{HTML}{0369A1}    % 조류 · 바람

% ---- 선과 화살촉 -----------------------------------------------------------
\tikzset{
  sig/.style   = {ink, line width=0.5pt, -{Stealth[length=4pt, width=3pt]}},
  wire/.style  = {ink, line width=0.5pt},
  fb/.style    = {ink, line width=0.5pt, densely dashed,
                  -{Stealth[length=4pt, width=3pt]}},
  asig/.style  = {accent, line width=0.5pt, densely dashed,
                  -{Stealth[length=4pt, width=3pt]}},
  % 블록 — 흰 바탕, 직각, 검은 테두리. 색은 테두리에만 준다
  blk/.style   = {draw=ink, line width=0.55pt, fill=white,
                  minimum width=18mm, minimum height=8.5mm, inner sep=2pt, font=\small},
  gblk/.style  = {blk, draw=guid},
  cblk/.style  = {blk, draw=ctrl},
  tblk/.style  = {blk, draw=thr},
  pblk/.style  = {blk, draw=plant},
  nblk/.style  = {blk, draw=nav},
  % 합산점 · 분기점
  sum/.style   = {draw=ink, line width=0.55pt, fill=white, circle,
                  inner sep=0pt, minimum size=4.6mm, font=\scriptsize},
  dot/.style   = {ink, fill=ink, circle, inner sep=0pt, minimum size=1.5mm},
  % 글자
  sname/.style = {font=\small, inner sep=1.5pt},
  note/.style  = {font=\scriptsize, soft, inner sep=1.5pt},
  wnote/.style = {font=\scriptsize, warn, inner sep=1.5pt},
  anote/.style = {font=\scriptsize, accent, inner sep=1.5pt},
  axis/.style  = {ink, line width=0.5pt},
  tick/.style  = {font=\scriptsize, soft},
  % 묶음 상자 — 서브시스템 표시
  grp/.style   = {draw=soft, line width=0.4pt, densely dotted, rounded corners=1mm,
                  inner sep=3mm},
}


\begin{document}
\begin{tikzpicture}[x=1mm, y=1mm]

% ===========================================================================
% 왼쪽 — 두 성분은 직각이다
% ===========================================================================
\begin{scope}[shift={(-52,0)}]

  \node[sname, anchor=south] at (-5,34)
       {두 성분은 서로 \textbf{직각} — 그래서 직각삼각형};

  \draw[axis, -{Stealth[length=3.5pt,width=2.8pt]}] (-24,0) -- (14,0);
  \draw[axis, -{Stealth[length=3.5pt,width=2.8pt]}] (0,-6) -- (0,29);
  \node[note, anchor=north] at (13,-1) {$\hat{e}_r$ 쪽};
  \node[note, anchor=east]  at (-1,28) {$\hat{e}_\theta$ 쪽};

  % 밑변 — rdot (음수라 왼쪽으로)
  \draw[guid, line width=0.9pt, -{Stealth[length=4pt,width=3pt]}] (0,0) -- (-15.5,0);
  \node[note, guid, anchor=north] at (-8,-1.4) {밑변 $=\dot{r}=-0.86$};

  % 높이 — r*thetadot (밑변 끝에서 올린다)
  \draw[thr, line width=0.9pt, -{Stealth[length=4pt,width=3pt]}] (-15.5,0) -- (-15.5,22.1);
  \node[note, thr, anchor=east, align=right] at (-17.2,13)
       {높이 $=r\dot{\theta}$\\$=+1.23$};

  % 직각 표시
  \draw[soft, line width=0.4pt] (-13.3,0) -- (-13.3,2.2) -- (-15.5,2.2);

  % 빗변 — v_d
  \draw[plant, line width=1.3pt, -{Stealth[length=5pt,width=4pt]}] (0,0) -- (-15.5,22.1);
  \node[note, plant, anchor=south west, align=left] at (-7.0,16.0)
       {빗변 $=v_d$\\$=1.50$};

  % alpha 호
  \draw[thr, line width=0.6pt, -{Stealth[length=3pt,width=2.4pt]}] (8,0) arc (0:124.9:8);
  \node[sname, thr, fill=white, inner sep=1pt] at (4.6,9.0) {$\alpha$};

  \node[note, anchor=north west, text width=56mm] at (-24,-9)
       {$\tan\alpha \;=\; \dfrac{\text{높이}}{\text{밑변}}
         \;=\; \dfrac{r\dot{\theta}}{\dot{r}}
         \;=\; \dfrac{1.23}{-0.86} \;=\; -1.43$};

  \node[note, anchor=north west, text width=46mm] at (-24,-24)
       {밑변이 \textbf{음수}라 화살표가 왼쪽 위를 향한다.
        그래서 $\alpha$ 는 $90^\circ$ 를 넘는다 ($124.9^\circ$)};

\end{scope}

% ===========================================================================
% 오른쪽 — 비 하나로는 각이 정해지지 않는다
% ===========================================================================
\begin{scope}[shift={(34,0)}]

  \node[sname, anchor=south] at (0,34)
       {비 $-1.43$ 을 주는 점은 \textbf{둘}이다};

  \draw[axis, -{Stealth[length=3.5pt,width=2.8pt]}] (-26,0) -- (26,0);
  \draw[axis, -{Stealth[length=3.5pt,width=2.8pt]}] (0,-28) -- (0,29);
  \node[note, anchor=north] at (25,-1) {$\dot{r}$};
  \node[note, anchor=east]  at (-1,28) {$r\dot{\theta}$};

  % 같은 기울기 위에 두 점이 있다
  \draw[soft, line width=0.4pt, densely dashed] (18.5,-26.4) -- (-18.5,26.4);

  % 점 A — 실제로 쓰는 쪽
  \draw[plant, line width=1.2pt, -{Stealth[length=5pt,width=4pt]}] (0,0) -- (-15.5,22.1);
  \fill[plant] (-15.5,22.1) circle (1.3);
  \draw[thr, line width=0.6pt, -{Stealth[length=3pt,width=2.4pt]}] (8,0) arc (0:124.9:8);
  \node[sname, thr, fill=white, inner sep=1pt] at (4.6,9.4) {$\alpha$};
  \node[note, plant, anchor=east, text width=29mm, align=right] at (-17,15)
       {$(\dot{r},\, r\dot{\theta})$\\
        $= (-0.86,\, +1.23)$\\
        $\alpha = +124.9^\circ$\\
        접근하며 시계로 돎};

  % 점 B — 비는 같지만 정반대
  \draw[errc, line width=1.2pt, -{Stealth[length=5pt,width=4pt]}] (0,0) -- (15.5,-22.1);
  \fill[errc] (15.5,-22.1) circle (1.3);
  \draw[errc, line width=0.6pt, -{Stealth[length=3pt,width=2.4pt]}] (13,0) arc (0:-55.1:13);
  \node[sname, errc, fill=white, inner sep=1pt] at (13.0,-5.2) {$\alpha$};
  \node[note, errc, anchor=west, text width=29mm] at (17,-15)
       {$(\dot{r},\, r\dot{\theta})$\\
        $= (+0.86,\, -1.23)$\\
        $\alpha = -55.1^\circ$\\
        멀어지며 반시계로 돎};

  \node[wnote, anchor=north, text width=60mm, align=center] at (0,-34)
       {두 침로는 정확히 $180^\circ$ 어긋나 있다.
        $\arctan(-1.43)$ 은 둘 다 $-55.1^\circ$ 로 답한다};

\end{scope}

% ===========================================================================
% 아래 — 읽는 법
% ===========================================================================
\node[sname, anchor=north west] at (-76,-52) {$\operatorname{atan2}$ 를 읽는 법};

\node[note, anchor=north west, text width=66mm] at (-76,-58)
     {$\operatorname{atan2}(y,\ x)$ 는 \textbf{원점에서 점 $(x,\ y)$ 를 본 각}이다.
      가로축에서 재고, 범위는 $-180^\circ$ 부터 $180^\circ$ 까지 한 바퀴 전부다.\\[2pt]
      인수 순서가 \textbf{(세로, 가로)} 로 뒤집혀 있다 — 헷갈리는 자리가 여기다.
      가로가 $\dot{r}$, 세로가 $r\dot{\theta}$ 이므로
      $\alpha = \operatorname{atan2}(r\dot{\theta},\ \dot{r})$ 로 쓴다};

\node[note, anchor=north west, text width=66mm] at (-4,-58)
     {$\arctan$ 으로 쓸 수 없는 까닭 셋\\[2pt]
      ① 두 성분의 \textbf{부호}를 잃는다 — 위 오른쪽 그림의 두 점을 구별하지 못한다\\[1pt]
      ② 값의 범위가 $-90^\circ \sim 90^\circ$ 뿐이다. 원을 도는 유도는 한 바퀴가 다 필요하다\\[1pt]
      ③ $r = r_d$ 일 때 $\dot{r} = 0$ 이라 나눗셈이 터진다.
      $\operatorname{atan2}$ 는 그 자리에서 $\pm 90^\circ$ 를 그대로 돌려준다};

\end{tikzpicture}
\end{document}
